% TOP_PROBLEM: {"id": 3000015, "title": "Complexity of computing the rotor-router action", "category_id": 2, "category": {"id": 2, "name": "combinatorics", "display_name": "Combinatorics", "description": "Counting problems, graph theory, discrete structures.", "slug": "combinatorics", "order_index": 2, "created_at": "2026-07-31T15:26:25.671Z"}, "status": "open"}
% FIRST_POSTED: null
% ENTRY_KIND: statement_only
% Imported from: batch09.tex; UnsolvedMath ID 3000015
\documentclass[11pt]{article}
\usepackage[margin=25mm]{geometry}
\usepackage{fontspec}
\setmainfont{Latin Modern Roman}
\usepackage{amsmath,amssymb,mathtools}
\usepackage{unicode-math}
\setmathfont{Latin Modern Math}
\usepackage{xurl,hyperref}
\hypersetup{colorlinks=true,urlcolor=blue}
\setcounter{secnumdepth}{0}
\setlength{\parindent}{0pt}
\setlength{\parskip}{5pt}
\setlength{\emergencystretch}{3em}
\newcommand{\sref}[2]{\hyperlink{source:#1:#2}{[S#2]}}
\newcommand{\cataloguescope}[1]{\paragraph{Recorded target.}#1}
\begin{document}
% UnsolvedMath ID 3000015; source catalogue code AMR-029-0015
\setcounter{section}{3000014}
\section{Complexity of computing the rotor-router action}\label{3000015}
{\small\sffamily UnsolvedMath ID 3000015 · Combinatorics}
\subsection{Definitions and mathematical statement}

\paragraph{Shared notation.}$\mathbb N=\{1,2,\ldots\}$, and $\mathbb Z,\mathbb Q,\mathbb R,\mathbb C$ denote the integers, rationals, reals and complex numbers. Any other conventions are specified locally.


Let $G$ be a finite connected undirected loopless multigraph. Supply a cyclic order of incident edges at every vertex (a ribbon structure) and a root $r$. Let $\Delta$ be its integer Laplacian, with diagonal degrees and off-diagonal entries minus edge multiplicities. Its sandpile group is $\operatorname{Pic}^0(G)=\{D\in\mathbb Z^{V(G)}:\sum_v D(v)=0\}/\Delta\mathbb Z^{V(G)}$. Orient a spanning tree $T$ toward $r$, giving one outgoing rotor at each other vertex. The action of $v-r$ on $T$ places one chip at $v$, repeatedly advances its current vertex rotor to the next edge in the cyclic order and moves the chip along that edge, stopping upon first reaching $r$. The resulting rotors form another tree. These commuting bijections and their inverses define the action of every degree-zero divisor, which factors through $\operatorname{Pic}^0(G)$. Encode graph adjacency multiplicities in binary. At each vertex encode the cyclic ribbon order as a cyclic list of runs of consecutive parallel arcs, recording the head of each run and its positive binary-encoded length. A rotor or tree arc is specified by its run and binary position within that run; an ordinary explicit edge list is the special case of runs of length one. The input consists of this graph and ribbon encoding, $r,T$ and a degree-zero divisor representative $D$ with binary-encoded integer coefficients. Complexity is measured in this compressed input length, so it must not assume time polynomial in the expanded number of parallel arcs.
\paragraph{Statement recorded in the supplied source.}
Determine the computational complexity of producing the spanning tree $D_r(T)$ obtained by this rotor-router sandpile-group action, for arbitrary such inputs. In particular, is there an algorithm polynomial in the input length rather than in the length of its individual routing simulation? \sref{3000015}{2}
\subsection{Short English statement}
Determine how efficiently a specified sandpile-group element can act on a spanning tree by rotor routing.
\subsection{Sources}
\begin{description}
\item[\hypertarget{source:3000015:1}{S1}] ULAM.AI and the UnsolvedMath Contributors, \emph{UnsolvedMath: A Curated Collection of Open Mathematics Problems}, record 3000015 (AMR-029-0015). CC BY 4.0. \url{https://huggingface.co/datasets/ulamai/UnsolvedMath}.
\item[\hypertarget{source:3000015:2}{S2}] L. Tóthmérész, Rotor-routing reachability is easy, chip-firing reachability is hard (2022), Section 3, definitions, and final paragraph following Corollary 3.9, pp.8 and15. \url{https://arxiv.org/pdf/2102.11970}.
\item[\hypertarget{source:3000015:3}{S3}] M. Chan, T. Church and J. A. Grochow, Rotor-routing and spanning trees on planar graphs (2014), Section 2, Definitions and notation, pp.3--4. \url{https://arxiv.org/pdf/1308.2677}.
\end{description}
\label{end:3000015}
\end{document}
